\documentclass{article}

\usepackage{amsmath} 
\usepackage{float}
\usepackage{graphicx}

\title{PHY2650 HW3}
\author{Chen Yilun(123090048)}



\begin{document}

\maketitle


\section{Q1}
\subsection{(a)}

\textbf{Results:}
\begin{tabular}{rl}
0 & : 0.0, \\
0.1 & : 0.09966766429033642, \\
0.2 & : 0.19736503092637106, \\
0.3 & : 0.29123788265696554, \\
0.4 & : 0.3796528397004763, \\
0.5 & : 0.46128100641279307, \\
0.6 & : 0.5351535268080815, \\
0.7 & : 0.6006856680827494, \\
0.8 & : 0.6576698563284026, \\
0.9 & : 0.7062415149635489, \\
1.0 & : 0.746824132812436, \\
1.1 & : 0.7800614325765766, \\
1.2 & : 0.8067447581231203, \\
1.3 & : 0.8277429893931986, \\
1.4 & : 0.8439407138416479, \\
1.5 & : 0.8561883936248748, \\
1.6 & : 0.8652662259422558, \\
1.7 & : 0.871861593406611, \\
1.8 & : 0.8765586342591011, \\
1.9 & : 0.8798376097200239, \\
2.0 & : 0.8820813907623578, \\
2.1 & : 0.8835864418783616, \\
2.2 & : 0.8845760209054464, \\
2.3 & : 0.8852138115715777, \\
2.4 & : 0.8856167458989558, \\
2.5 & : 0.88586627361749, \\
2.6 & : 0.8860177453974631, \\
2.7 & : 0.8861078761616262, \\
2.8 & : 0.886160446739861, \\
2.9 & : 0.886190503406593, \\
3.0 & : 0.8862073482595122
\end{tabular}

\subsection{(b)Plot}
\begin{figure}[H]
    \centering
    \includegraphics[width=0.75\linewidth]{2650HW3/HW3Q1b.png}
    \caption{Plot of E(x) vs x}
    \label{fig:enter-label}
\end{figure}


\section{Q2}
\subsection{Results:}
Estimated error: 0.026633333333333137\\
Real error: 0.10655999999999999


\subsection{Analysis}

$I_{1}$ with 10 slices gives $E(1) = 4.50656$ which is not a good estimation for the real value 4.4, and it may be due to the insufficient number of slices. Hence the true error does not agree perfectly with the estimated one when using the value of $I_{1}$.


\section{Q3}
\subsection{Output}
\textbf{Results: To achieve the accuracy of $\epsilon < 10 ^ {-4}$ }

(1)Number of intervals for \textbf{left-hand rectangle rule}: 4194304

\textbf{(Estimateing the error by comparing with the result 
from the trapezoidal rule)}

Integral using \textbf{left-hand rectangle rule}: 0.5672623414436723

(2)Number of intervals for \textbf{adaptive trapezoidal rule}: 2097152

Integral using adaptive \textbf{adaptive trapezoidal rule}: 0.5672622052304107

(3)Number of intervals for \textbf{adaptive Simpson's rule}: 65536

Integral using \textbf{adaptive Simpson's rule}: 0.564498512178105

(4)As for the \textbf{Romberg method:}

Iteration:     1, I =   -0.267060634995734, error estimate = 0.304259196878092

Iteration:     2, I =   -0.041404061118524, error estimate = 0.301721373096733

Iteration:     2, I =   -0.026360289526710, error estimate = 0.240700345469024

Iteration:     3, I =    0.130079387755442, error estimate = 0.246913792148150

Iteration:     3, I =    0.141511617680373, error estimate = 0.182915678798897

......

Iteration:    24, I =    0.567679056654760, error estimate = 0.000122695129418

Iteration:    25, I =    0.567757971398380, error estimate = 0.000139234347817

Iteration:    25, I =    0.567763971903305, error estimate = 0.000096008078806

There are totally 25 Iterations, $2^{25}$ slices.

Integral using \textbf{Romberg method:} 0.5677639719033053

\subsection{Plot}

\begin{figure}[H]
    \centering
    \includegraphics[width=0.7\textwidth]{2650HW3/Q3LFRRRR.png} % 需要替换为实际的图片文件名
    \caption{}
\end{figure}

\begin{figure}[H]
    \centering
    \includegraphics[width=0.7\textwidth]{2650HW3/Q3TPZ.png} % 需要替换为实际的图片文件名
    \caption{}
\end{figure}

\begin{figure}[H]
    \centering
    \includegraphics[width=0.7\textwidth]{2650HW3/Q3SPS.png} % 需要替换为实际的图片文件名
    \caption{}
\end{figure}

\begin{figure}[H]
    \centering
    \includegraphics[width=0.7\textwidth]{2650HW3/Q3Romberg.png} % 需要替换为实际的图片文件名
    \caption{}
\end{figure}

\end{document}


